\clearpage
\section{Recommendations for migrating IPsec VPNs to post-quantum cryptography}
\label{sec:reco}

% CONFIDENTIALITY: this chapter is public. It is based only on the lab results
% of this thesis and on public regulation. No information about the employer
% (audit findings, architecture, vendors, inventory) goes here unless approved
% in writing. The mapping onto the bank's environment is a separate internal
% document.

\placeholder{Recommendations for operators of IPsec VPNs in regulated financial
institutions and IoT deployments, derived from the results in
Section~\ref{sec:results}.}

\subsection{Regulatory and organisational context}
\label{sec:reco-context}
\placeholder{Obligations under the Act on the National Cybersecurity System
(UKSC) relevant to cryptography and risk management of essential services in the
banking sector; relation to NIS2 and DORA; EU PQC roadmap milestones
(see Section~\ref{sec:sota-regulatory}).}

\subsection{Cryptographic inventory and prioritisation}
\label{sec:reco-inventory}
\placeholder{How to inventory IPsec tunnels and their algorithms; prioritising
tunnels by required confidentiality lifetime of the data (HNDL exposure).}

\subsection{Configuration recommendations}
\label{sec:reco-config}
\placeholder{Recommended hybrid combinations and parameter sets, rekeying policy
given the measured cost, IKE fragmentation settings, PPK as an interim measure
where ML-KEM is not yet available.}

\subsection{Network and vendor readiness}
\label{sec:reco-readiness}
\placeholder{Vendor support for RFC 9370, RFC 9242, RFC 7383 and ML-KEM;
handling of fragmented UDP 500/4500 by firewalls and NAT devices;
constraints of IoT gateways.}

\subsection{Migration plan and verification}
\label{sec:reco-plan}
\placeholder{Phased rollout, interoperability testing, monitoring of handshake
failures and fragment loss, fallback strategy.}
